\section{Tokenization Analysis}
\label{app:tokenization}

\begin{figure*}[t!]
        \centering      
        \begin{subfigure}[b]{0.3\textwidth}  
            \includegraphics[width=\textwidth]{figures/log_counts.png}
            \caption[]%
            {{\small Count analysis\vspace{0.25cm}}}    
            \label{fig:tokenization-frequency-analysis}
        \end{subfigure} 
        \begin{subfigure}[b]{0.3\textwidth}   
            \includegraphics[width=\textwidth]{figures/fertility.png}
            \caption[]%
            {{\small Fertility analysis}}   
            \label{fig:tokenization-fertility-analysis}
        \end{subfigure}
        \begin{subfigure}[b]{0.3\textwidth}   
            \includegraphics[width=\textwidth]{figures/whitespace.png}
            \caption[]%
            {{\small Whitespace analysis}}   
            \label{fig:tokenization-whitespace-analysis}
        \end{subfigure}
        \caption[]{Tokenization analysis. Tokens with small IDs, which have a high count in the tokenizer training data, also tend to have a high count in \dolma (a). The Stack has a substantially higher fertility compared to the other data sources (b), which can be explained by the higher relative frequency of whitespace characters such ``\textbackslash n'' and ``\textbackslash t'' (c). See text for more details.}
        \label{fig:tokenization-analysis}
\end{figure*}



The first step of processing text with LMs is \emph{tokenization}, i.e., mapping the text to a sequence of tokens with corresponding input embeddings \citep{sennrich-etal-2016-neural, kudo-2018-subword, kudo-richardson-2018-sentencepiece}. Recently, there has been a growing interest in the question of how well LM tokenizers fit different data sources \citep[e.g., data in different languages;][]{ahia2023languages, petrov2023language}
Inspired by this emerging line of work, we conduct an explorative analysis of the GPTNeoX tokenizer~\citep{Black2022GPTNeoX20BAO} applied to \dolma, which provides a first picture of how challenging the different data sources comprised by \dolma are for current LM tokenizers. 

We start by taking a global look at the tokenizer's fit to \dolma. Out of the 50,280 tokens in the tokenizer vocabulary, 50,057 are present in the tokenized text of \dolma. In other words, 223 tokens are never used, amounting to roughly 0.4\% of the tokenizer vocabulary. The 223 tokens mostly consist of combinations of whitespace characters (e.g., ``\textbackslash n\textbackslash n {} '', two newline characters followed by two blank space characters). Note that when training an LM with the examined tokenizer on \dolma, the input embeddings corresponding to these tokens would not be updated. 
In terms of the count distribution of tokens, we find that tokens with smaller IDs tend to have higher counts in \dolma (see Figure \ref{fig:tokenization-frequency-analysis}), which is also reflected by a strong Spearman's correlation between (i) the ranking of tokens based on their counts in \dolma and (ii) the token IDs ($r =$ 0.638, $p <$ 0.001). Given how the tokenizer was trained \citep{sennrich-etal-2016-neural, Black2022GPTNeoX20BAO}, smaller IDs correspond to byte pairs merged earlier and hence tokens occurring more frequently in the tokenizer training data 
Overall, these results suggest a good fit of the GPTNeoX tokenizer to \dolma. 


Does the tokenizer fit all data sources included in \dolma equally well? To examine this question, we analyze fertility, which is defined as the average number of tokens per word generated by a tokenizer \citep{acs2019,Scao2022BLOOMA1}, in our case measured on a specific data source.
We find that fertility is similar for most data sources, ranging between 1.15 (conversational forum subset) and 1.28 (books subset), with the exception of the code subset, which has a substantially higher fertility of 2.45 (see Figure \ref{fig:tokenization-fertility-analysis}). This means that the costs of processing the code subset --- be they computational or financial in nature \citep{petrov2023language} --- are more than twice as high compared to the other data sources. 

What causes this discrepancy? We find that in the code subset (which mostly contains code), words are often preceded by whitespace characters \emph{other than} 
a blank space (e.g., newline, tab, return). Crucially, while a blank space before a word is tokenized as part of that word (e.g., \textit{I love you} $\rightarrow$ ``I'', `` love'', `` you''), other whitespace characters yield separate tokens (e.g., \textit{I \quad love \quad you} $\rightarrow$ ``I'', ``\textbackslash t'', ``love'', ``\textbackslash t'', ``you''). This can also be seen by plotting the relative frequency of tokens representing whitespace characters by data source, which is one order of magnitude higher for The Stack compared to most other data sources (see Figure \ref{fig:tokenization-whitespace-analysis}). When training LMs on The Stack (or code more generally), it thus might be advisable to add special tokens to the tokenizer \citep[e.g., ``\textbackslash nif'';][]{hong-etal-2021-avocado}. It is important to notice that this observation applies to most tokenizers in use today (e.g., the tokenizer used by GPT-4), which tend to lack tokens such as ``\textbackslash nif''.

